\documentclass[11pt]{article}

\usepackage[margin=1in]{geometry}
\usepackage{amsmath,amssymb}
\usepackage{booktabs}
\usepackage{longtable}
\usepackage{array}
\usepackage{tabularx}
\usepackage{graphicx}
\usepackage{enumitem}
\usepackage[hidelinks]{hyperref}
\usepackage{caption}
\usepackage{float}

\captionsetup{font=small,labelfont=bf}
\setlist[itemize]{leftmargin=1.5em,itemsep=0.25em,topsep=0.25em}

\newcommand{\Var}{\operatorname{Var}}
\newcommand{\NA}{\multicolumn{1}{c}{--}}

\title{Empirical Analysis Section Extracted from FT-PPI}
\author{}
\date{}

\begin{document}
\maketitle

\setcounter{section}{5}
\section{Empirical Analysis}

In this section, we present empirical evidence on the Wine Reviews dataset to validate the
proposed fine-tuning and prediction-powered inference framework. The empirical analysis is
organized around the two ingredients of the method. First, we examine whether the residual
variance of a fine-tuned LLM follows the scaling law that underlies the label-allocation rule.
Second, we evaluate whether this allocation principle leads to improved finite-population mean
estimation when only a small labeled dataset is available.

To keep these two questions statistically separate, we use an independent pilot set only for
scaling-law validation and reserve a disjoint target population for all inference experiments.
Section~6.1 describes the dataset, estimand, label accounting, and fine-tuning protocol.
Section~6.2 validates the scaling-law premise. Section~6.3 tests the full-grid allocation
objective on fresh target-population splits. Section~6.4 studies whether a deployable ramp-up
procedure can recover a low-regret allocation without training on the full grid. Section~6.5
compares the resulting FT+PPI estimators with natural benchmarks. Section~6.6 discusses how the
same allocation rule may be combined with Cross-PPI. Additional diagnostics for residual
variance, prediction stability, and label leakage are reported in Appendix~\ref{app:diagnostics}.

\subsection{Empirical Setup}
\label{subsec:empirical_setup}

Textual product reviews are a common source of information for assessing perceived product
quality in many markets. Such reviews contain rich signals about aroma, flavor, balance, and
finish, but they are costly to aggregate when reliable expert ratings are needed. The empirical
task therefore mirrors a common managerial problem: a firm observes many pieces of product text
but can afford only a limited number of human labels, and it wants to estimate an aggregate
quality measure with valid uncertainty quantification.

\paragraph{Dataset and Estimand.}
We study this problem using the Wine Reviews dataset, which contains reviews written by
professional wine critics at \emph{Wine Enthusiast} magazine and publicly available on Kaggle
(Thoutt 2017). This dataset has also been used in prior work on prediction-powered inference and
surrogate outcomes (Ji et al. 2025). Each observation contains a free-form review text and an
expert rating from 80 to 100 points. We use only the review text as the model input and the
expert rating as the outcome.

After dropping observations with missing review text or rating and deduplicating by review
description, the cleaned dataset contains 119,955 unique reviews. Due to the computational cost
of repeated LLM fine-tuning, we randomly sample 25,000 reviews from this cleaned dataset and use
them as the experimental finite population
\[
    \mathcal{P}_0=\{(X_i,Y_i):i=1,\ldots,N_0\},\qquad N_0=25{,}000.
\]
The raw rating variance in $\mathcal{P}_0$ is approximately 9.57. This number is the residual
variance of the constant predictor that always predicts the population mean, and serves as a
baseline for judging whether fine-tuned LLM surrogates extract useful information from text.
Table~\ref{tab:wine_examples} reports two representative examples from the Wine Reviews dataset.

\begin{table}[H]
\centering
\caption{Examples from the Wine Reviews dataset}
\label{tab:wine_examples}
\begin{tabularx}{\textwidth}{@{}p{0.16\textwidth} X@{}}
\toprule
Review ($X$) &
This is a walk backward after the impressive 2012. Almost impenetrably black, the flavors
converge around espresso and bitter chocolate, yet the tannins have a green edge. The wine
simply feels flat in the mouth with no life to it. \\
\addlinespace
Rating ($Y$) & 85 \\
\midrule
Review ($X$) &
This is one of the great Rieslings from the Wachau, a wonderful panoply of ripe, tropical
fruit, pierced with flint, spice and minerality. It is rich and opulent, while never losing
sight of the core tautness of a fine Riesling. \\
\addlinespace
Rating ($Y$) & 95 \\
\bottomrule
\end{tabularx}
\end{table}

\paragraph{Pilot and target populations.}
Before running any target-budget experiment, we split the experimental population once into two
disjoint sets:
\[
    \mathcal{P}_0=\mathcal{H}_{\mathrm{scale}}\,\dot\cup\,\mathcal{P}_{\mathrm{target}}.
\]
The pilot set has size $|\mathcal{H}_{\mathrm{scale}}|=10{,}000$, and the target population has
size $|\mathcal{P}_{\mathrm{target}}|=15{,}000$. The set
$\mathcal{H}_{\mathrm{scale}}$ is used only in Section~6.2 to assess whether LoRA-Var residual
variance follows a stable power law. It is not used in the allocation validation, ramp-up
recovery, estimator comparison, or diagnostic leakage checks for the target population. All
inferential experiments draw labels only from $\mathcal{P}_{\mathrm{target}}$, whose
finite-population mean is the target estimand.

\paragraph{PPI Estimator Construction.}
In Sections~6.3--6.5, for each replication $r$ and labeled budget $B$, we draw a fresh small
labeled dataset
\[
    \mathcal{N}_r^{(B)}\subset\mathcal{P}_{\mathrm{target}},\qquad
    |\mathcal{N}_r^{(B)}|=B,
\]
by simple random sampling without replacement. The remaining reviews
\[
    \mathcal{M}_r^{(B)}
    =
    \mathcal{P}_{\mathrm{target}}\setminus\mathcal{N}_r^{(B)}
\]
form the large unlabeled text dataset. As in the main framework, the estimator may use the texts
$X_i$ for $i\in\mathcal{M}_r^{(B)}$, but the labels $Y_i$ for these reviews are never used for
training, model selection, allocation selection, scaling-law fitting, correction, or estimator
construction.

Within the labeled budget, we first split
\[
    \mathcal{N}_r^{(B)}
    =
    \mathcal{V}_r^{(B)}\,\dot\cup\,\mathcal{N}_{r,\mathrm{eff}}^{(B)} .
\]
The validation subset $\mathcal{V}_r^{(B)}$ is used for checkpoint selection and, when
applicable, for target-only ramp-up and allocation selection. The remaining labels
$\mathcal{N}_{r,\mathrm{eff}}^{(B)}
=\mathcal{N}_r^{(B)}\setminus\mathcal{V}_r^{(B)}$ are the only labels available for the final
estimator.
For a candidate allocation $s$, we further split
\[
    \mathcal{N}_{r,s}^{\mathrm{FT}}\subset \mathcal{N}_{r,\mathrm{eff}}^{(B)},
    \qquad |\mathcal{N}_{r,s}^{\mathrm{FT}}|=s,
    \qquad
    \mathcal{N}_{r,s}^{\mathrm{PPI}}
    =
    \mathcal{N}_{r,\mathrm{eff}}^{(B)}
    \setminus
    \mathcal{N}_{r,s}^{\mathrm{FT}}.
\]
The set $\mathcal{N}_{r,s}^{\mathrm{FT}}$ is used for fine-tuning, and
$\mathcal{N}_{r,s}^{\mathrm{PPI}}$ is used for PPI correction. Thus, validation, fine-tuning, and
correction labels are disjoint within each replication and budget. Table~\ref{tab:label_accounting}
summarizes the nested decomposition and role of each split.

\begin{table}[H]
\centering
\caption{Nested data decomposition in the empirical protocol}
\label{tab:label_accounting}
\footnotesize
\setlength{\tabcolsep}{5pt}
\renewcommand{\arraystretch}{1.13}
\begin{tabularx}{\textwidth}{@{}p{0.25\textwidth}p{0.37\textwidth}X@{}}
\toprule
Set & Composition / size & Role \\
\midrule
$\mathcal{P}_0$ &
$\mathcal{H}_{\mathrm{scale}}\,\dot\cup\,\mathcal{P}_{\mathrm{target}}$,
$|\mathcal{P}_0|=25{,}000$ &
Experimental finite population. \\
\addlinespace[0.15em]
$\quad\mathcal{H}_{\mathrm{scale}}$ &
$|\mathcal{H}_{\mathrm{scale}}|=10{,}000$ &
Used only for independent scaling-law validation in Section~6.2. \\
\addlinespace[0.15em]
$\quad\mathcal{P}_{\mathrm{target}}$ &
$\mathcal{P}_0\setminus\mathcal{H}_{\mathrm{scale}}$,
$|\mathcal{P}_{\mathrm{target}}|=15{,}000$ &
Target finite population for allocation, ramp-up, and inference experiments. \\
\midrule
$\quad\quad\mathcal{N}_r^{(B)}$ &
$\mathcal{N}_r^{(B)}\subset\mathcal{P}_{\mathrm{target}}$,
$|\mathcal{N}_r^{(B)}|=B$ &
Small labeled dataset drawn in replication $r$. \\
\addlinespace[0.15em]
$\quad\quad\mathcal{M}_r^{(B)}$ &
$\mathcal{P}_{\mathrm{target}}\setminus\mathcal{N}_r^{(B)}$ &
Large unlabeled text dataset; labels are never used. \\
\midrule
$\quad\quad\quad\mathcal{V}_r^{(B)}$ &
$\mathcal{V}_r^{(B)}\subset\mathcal{N}_r^{(B)}$
(20\% of $B$ in allocation experiments) &
Validation for checkpoint selection and allocation/ramp-up decisions. \\
\addlinespace[0.15em]
$\quad\quad\quad\mathcal{N}_{r,\mathrm{eff}}^{(B)}$ &
$\mathcal{N}_r^{(B)}\setminus\mathcal{V}_r^{(B)}$ &
Effective labels available to the final FT+PPI estimator. \\
\midrule
$\quad\quad\quad\quad\mathcal{N}_{r,s}^{\mathrm{FT}}$ &
$\mathcal{N}_{r,s}^{\mathrm{FT}}\subset
\mathcal{N}_{r,\mathrm{eff}}^{(B)}$,
$|\mathcal{N}_{r,s}^{\mathrm{FT}}|=s$ &
Labels used to fine-tune the surrogate. \\
\addlinespace[0.15em]
$\quad\quad\quad\quad\mathcal{N}_{r,s}^{\mathrm{PPI}}$ &
$\mathcal{N}_{r,\mathrm{eff}}^{(B)}
\setminus\mathcal{N}_{r,s}^{\mathrm{FT}}$ &
Correction labels for the PPI rectifier. \\
\bottomrule
\end{tabularx}
\end{table}

\paragraph{LLM Fine-Tuning.}
We use \texttt{Qwen/Qwen2.5-1.5B-Instruct} as the base LLM. Given a review text $X$, the
model forms a pooled representation from the last non-padding token, and a scalar linear head
maps this representation to the prediction $f(X)$. Ratings are standardized as
\[
    y_{\mathrm{scaled}}=(y_{\mathrm{raw}}-90)/5,
    \qquad
    \hat y_{\mathrm{raw}}=90+5\hat y_{\mathrm{scaled}},
\]
so that optimization is stable while all reported residual variances remain interpretable on the
original rating scale. We use LoRA adapters on the linear modules and train the scalar head
jointly with the adapters.

The standard objective for rating prediction is mean-squared error. In contrast, the downstream
PPI estimator is governed by the variance of the residual $Y-f(X)$ after bias correction. We
therefore also fine-tune with a residual-variance loss. For a mini-batch of size $k$, define
$r_i=Y_i-f(X_i)$ and $\bar r=k^{-1}\sum_{i=1}^k r_i$. The variance loss is
\[
    L_{\mathrm{var}}=\frac{1}{k}\sum_{i=1}^k(r_i-\bar r)^2.
\]
This objective is aligned with the PPI correction stage: the fine-tuning stage is used to reduce
residual variance, while systematic bias is removed by the correction labels in
$\mathcal{N}_{r,s}^{\mathrm{PPI}}$.

\subsection{Independent Scaling-Law Validation}
\label{subsec:independent_scaling_law}

Before using a scaling-law-based allocation rule, we first verify that the residual variance of
LoRA-Var fine-tuning is well described by a power law. This validation is deliberately separated
from the target-population experiments. It uses only the pilot set $\mathcal{H}_{\mathrm{scale}}$ and
asks whether
\[
    \Var\!\left(Y-f^{(s)}(X)\right)=a s^{-\alpha}+b
\]
is a reasonable empirical approximation over the range of fine-tuning sizes relevant for our
allocation experiments.

The pilot protocol separates model selection from scaling-law evaluation. In each pilot
replication, we randomly partition
\[
    \mathcal{H}_{\mathrm{scale}}
    =
    \mathcal{V}_{\mathrm{stop}}^{\mathrm{scale}}
    \,\dot\cup\,
    \mathcal{V}_{\mathrm{scale}}^{\mathrm{scale}}
    \,\dot\cup\,
    \mathcal{T}_{\mathrm{scale}},
\]
where
$|\mathcal{V}_{\mathrm{stop}}^{\mathrm{scale}}|=1{,}000$,
$|\mathcal{V}_{\mathrm{scale}}^{\mathrm{scale}}|=1{,}000$, and
$|\mathcal{T}_{\mathrm{scale}}|=8{,}000$. The set
$\mathcal{V}_{\mathrm{stop}}^{\mathrm{scale}}$ is used only for early stopping and checkpoint
selection. The set $\mathcal{V}_{\mathrm{scale}}^{\mathrm{scale}}$ is used only to measure
residual variance and fit the scaling law. The remaining set $\mathcal{T}_{\mathrm{scale}}$
contains the labels from which all fine-tuning subsets are drawn.

We run 10 independent pilot replications. In each replication, the fine-tuning subsets are nested
inside $\mathcal{T}_{\mathrm{scale}}$ over
\[
    s\in\{50,100,150,200,250,300,350,400,450,500,550,600,650,700\}.
\]
For each $s$, we train a LoRA-Var model on the first $s$ labels in the nested training order and
evaluate its residual variance on $\mathcal{V}_{\mathrm{scale}}^{\mathrm{scale}}$. We then fit the
power law to the mean residual variance across replications, reported on the original rating
scale. Error bars show 95\% confidence intervals.

\begin{figure}[H]
\centering
\begin{minipage}{0.47\textwidth}
\centering
\includegraphics[width=\linewidth]{artifacts/hyak/wine_loss_controlled_b10000_r10_stopmetric_s50_to700/figures/var_stop_var_validation_scale_residual_variance_power_law_linear_raw.pdf}
\vspace{0.35em}

(a) Power-law fit in linear space
\end{minipage}
\hfill
\begin{minipage}{0.47\textwidth}
\centering
\includegraphics[width=\linewidth]{artifacts/hyak/wine_loss_controlled_b10000_r10_stopmetric_s50_to700/figures/var_stop_var_validation_scale_residual_variance_power_law_loglog_raw.pdf}
\vspace{0.35em}

(b) Linear fit in log-log space
\end{minipage}
\caption{Independent validation of the variance-based scaling law on $\mathcal{H}_{\mathrm{scale}}$.
The fitted parameters are $\hat a=17.40$, $\hat\alpha=0.574$, and $\hat b=2.827$ on the
original rating scale.}
\label{fig:scaling_law_validation}
\end{figure}

Figure~\ref{fig:scaling_law_validation} shows a strong scaling-law fit. The fitted curve has
$R^2=0.995$, and the mean residual variance decreases from 4.67 at $s=50$ to 3.23 at
$s=700$, far below the constant-predictor baseline of 9.57. The log-log panel gives a second
view of the same pattern: after subtracting the fitted noise floor, the residual variance is close
to linear in $\log(s)$, with slope $-\hat\alpha=-0.574$.

This result validates the premise that additional fine-tuning labels reduce residual variance in a
smooth and predictable way. It does not, however, determine the final allocation in the target
population. The next two experiments evaluate allocation directly using fresh labeled budgets from
$\mathcal{P}_{\mathrm{target}}$. Additional implementation details for the scaling-law fit,
including the early-stopping split and diagnostic quantities, are given in
Appendix~\ref{app:scaling_details}.

\subsection{Full-Grid Allocation Validation}
\label{subsec:full_grid_allocation}

We next test whether the allocation rule implied by the scaling law identifies a useful
fine-tuning fraction on the target population. This experiment differs from Section~6.2 in its
role. Section~6.2 estimates the scaling-law parameters on the independent pilot set
$\mathcal{H}_{\mathrm{scale}}$; Section~6.3 keeps those parameters fixed and asks whether the
resulting theoretical allocation agrees with a full-grid empirical search on fresh labeled
budgets from $\mathcal{P}_{\mathrm{target}}$. To make the comparison interpretable, we run the
same full-grid protocol for LoRA-MSE + PPI as an empirical benchmark, but the theoretical
allocation itself is computed from the LoRA-Var scaling law.

For each budget $B$ and replication $r$, we draw a fresh small labeled dataset
$\mathcal{N}_r^{(B)}$ from $\mathcal{P}_{\mathrm{target}}$. The validation set is charged to this
same budget: we set
\[
    |\mathcal{V}_r^{(B)}|=0.2B,
    \qquad
    n_B
    =
    B-|\mathcal{V}_r^{(B)}|
    =
    0.8B .
\]
Only the remaining $n_B$ labels,
$\mathcal{N}_{r,\mathrm{eff}}^{(B)}
=\mathcal{N}_r^{(B)}\setminus\mathcal{V}_r^{(B)}$, are available for the final FT+PPI estimator.
Thus, an allocation fraction $\rho$ corresponds to
\[
    s(\rho)=\mathrm{round}(\rho n_B)
    \quad\text{fine-tuning labels, and}\quad
    n_B-s(\rho)
    \quad\text{PPI correction labels}.
\]
This accounting is important: the validation labels are used for allocation validation and are
not returned to the final fine-tuning or correction sets.

The theoretical allocation uses the scaling-law parameters estimated in Section~6.2. Plugging
$(\hat a,\hat\alpha,\hat b)$ into the FT-PPI allocation objective gives
\[
    \rho_B^{\mathrm{th}}
    =
    \arg\min_{0<\rho<1}
    \frac{
        \hat a\,s(\rho)^{-\hat\alpha}+\hat b
    }{
        n_B-s(\rho)
    },
    \qquad s(\rho)=\mathrm{round}(\rho n_B).
\label{eq:target_theoretical_allocation}
\]
Equivalently, the numerator predicts the residual variance of the LoRA-Var surrogate trained on
$s(\rho)$ labels, while the denominator is the number of labels left for PPI correction. We solve
this one-dimensional problem numerically for each $B$ after removing the validation labels from
the budget. This theoretical rule targets the dominant label-allocation term in equation (2);
the empirical full-grid diagnostic below evaluates the full plug-in variance, including the
unlabeled prediction-variance term.

To validate this theoretical optimum, we also run a full-grid empirical search over
\[
    \mathcal{G}
    =
    \{0.025,0.05,0.075,0.10,0.125,0.15,0.20,0.30,0.40,0.50,0.70,0.90\}.
\]
In addition, for LoRA-Var + PPI we train the exact theoretical allocation
$\rho_B^{\mathrm{th}}$ as an additional candidate. This candidate uses the same replication,
the same labeled budget, the same validation set, and the same nested ordering of
$\mathcal{N}_{r,\mathrm{eff}}^{(B)}$ as the grid candidates; only the number of fine-tuning
labels changes. Thus the LoRA-Var candidate set is
\[
    \mathcal{G}_B^{\mathrm{Var}}
    =
    \mathcal{G}\cup\{\rho_B^{\mathrm{th}}\},
\]
whereas LoRA-MSE + PPI is evaluated on $\mathcal{G}$ only because we do not propose a
separate MSE allocation rule. For each candidate $\rho$ and each fine-tuning objective
$m\in\{\mathrm{MSE},\mathrm{Var}\}$, we fine-tune a LoRA predictor on
$\mathcal{N}_{r,s(\rho)}^{\mathrm{FT}}$ and use the remaining
$\mathcal{N}_{r,s(\rho)}^{\mathrm{PPI}}$ labels for PPI correction. The main allocation metric is
the plug-in estimate of the FT+PPI estimator variance from equation (2):
\[
    \widehat V_{r,B}^{m}(\rho)
    =
    \frac{
        \widehat{\Var}_{\mathcal{N}_{r,s(\rho)}^{\mathrm{PPI}}}
        \{Y-f_{r,\rho}^{m}(X)\}
    }{
        n_B-s(\rho)
    }
    +
    \frac{
        \widehat{\Var}_{\mathcal{M}_r^{(B)}}\{f_{r,\rho}^{m}(X)\}
    }{
        |\mathcal{M}_r^{(B)}|
    } .
    \label{eq:full_grid_plugin_variance}
\]
Here $\widehat{\Var}_{A}$ denotes the sample variance over the set $A$, and
$f_{r,\rho}^{m}$ denotes the surrogate trained with loss $m$ at allocation $\rho$ in replication
$r$. The first term is the variance contribution from the PPI correction labels, and the second
term is the variance contribution from averaging predictions over the unlabeled text pool. This
quantity is the empirical counterpart of the allocation objective in the FT-PPI theory.

We aggregate the objective across the $R=10$ target-population replications:
\[
    \overline V_B^{m}(\rho)
    =
    \frac{1}{R}\sum_{r=1}^R
    \widehat V_{r,B}^{m}(\rho).
\]
The full-grid empirical benchmark for method $m$ is
\[
    \rho_B^{\mathrm{grid},m}
    =
    \arg\min_{\rho\in\mathcal{G}_B^m}
    \overline V_B^{m}(\rho).
\]
Here $\mathcal{G}_B^{\mathrm{Var}}=\mathcal{G}\cup\{\rho_B^{\mathrm{th}}\}$ and
$\mathcal{G}_B^{\mathrm{MSE}}=\mathcal{G}$. This convention makes the empirical benchmark fair:
if the theoretical LoRA-Var allocation is actually the best trained candidate, it is allowed to
win the benchmark rather than being compared only with nearby grid points.
For LoRA-Var + PPI, the performance columns are evaluated at the theoretical allocation
$\rho_B^{\mathrm{th}}$, and the grid-best allocation is reported only as the empirical benchmark.
For LoRA-MSE + PPI, there is no studied allocation rule; its performance is evaluated at
$\rho_B^{\mathrm{grid},\mathrm{MSE}}$ as a strong full-grid comparator.

We report two additional quantities to make the diagnostic interpretable. First, the normalized
variance ratio compares the estimated variance at the reported allocation to the estimated
variance of the Sample Mean under the same budget:
\[
    \mathrm{VarRatio}_B^m
    =
    \frac{
        \overline V_B^m(\rho_B^{\mathrm{eval},m})
    }{
        \overline V_B^{\mathrm{SM}}
    },
    \qquad
    \rho_B^{\mathrm{eval},\mathrm{Var}}=\rho_B^{\mathrm{th}},
    \quad
    \rho_B^{\mathrm{eval},\mathrm{MSE}}=\rho_B^{\mathrm{grid},\mathrm{MSE}}.
\]
Second, for LoRA-Var + PPI, allocation regret measures the loss from using the theoretical
allocation instead of the full-grid empirical benchmark:
\[
    \mathrm{Regret}_B^{\mathrm{Var}}
    =
    \frac{
        \overline V_B^{\mathrm{Var}}(\rho_B^{\mathrm{th}})
        -
        \overline V_B^{\mathrm{Var}}(\rho_B^{\mathrm{grid},\mathrm{Var}})
    }{
        \overline V_B^{\mathrm{Var}}(\rho_B^{\mathrm{grid},\mathrm{Var}})
    } .
\]
Finally, each replication also produces a PPI standard error
$\widehat{\mathrm{SE}}_{\mathrm{PPI},r,B}^{m}(\rho)$ and a within-replication inference interval
\[
    \mathrm{CI}_{\mathrm{PPI},r,B}^{m}(\rho)
    =
    \hat\mu_{r,B}^{m}(\rho)
    \pm
    1.96\,\widehat{\mathrm{SE}}_{\mathrm{PPI},r,B}^{m}(\rho).
\]
The coverage column is the empirical coverage of these within-replication PPI intervals at the
reported allocation:
\[
    \mathrm{Coverage}_{B}^{m}(\rho)
    =
    \frac{1}{R}
    \sum_{r=1}^R
    \mathbf{1}\!\left\{
    \mu_{\mathrm{target}}\in
    \mathrm{CI}_{\mathrm{PPI},r,B}^{m}(\rho)
    \right\}.
\]
Coverage is not used to choose the allocation; it checks that shorter variance estimates do not
come from invalid confidence intervals.

\begin{table}[H]
\centering
\caption{Full-grid allocation validation using the FT+PPI variance objective}
\label{tab:full_grid_allocation}
\scriptsize
\setlength{\tabcolsep}{2.7pt}
\renewcommand{\arraystretch}{1.08}
\begin{tabular}{@{}c c c l c c c c c c@{}}
\toprule
$B$ & $|\mathcal{V}|$ & $n_B$ & Method & Theory $\rho$ & Emp.-best $\rho$ & Eval. $\rho$ & Var ratio & Regret & PPI cov. \\
\midrule
300 & 60  & 240 & LoRA-MSE + PPI & \NA & \NA & \NA & \NA & \NA & \NA \\
    & 60  & 240 & LoRA-Var + PPI & \NA & \NA & \NA & \NA & \NA & \NA \\
\midrule
500 & 100 & 400 & LoRA-MSE + PPI & \NA & \NA & \NA & \NA & \NA & \NA \\
    & 100 & 400 & LoRA-Var + PPI & \NA & \NA & \NA & \NA & \NA & \NA \\
\midrule
700 & 140 & 560 & LoRA-MSE + PPI & \NA & \NA & \NA & \NA & \NA & \NA \\
    & 140 & 560 & LoRA-Var + PPI & \NA & \NA & \NA & \NA & \NA & \NA \\
\midrule
1000 & 200 & 800 & LoRA-MSE + PPI & \NA & \NA & \NA & \NA & \NA & \NA \\
     & 200 & 800 & LoRA-Var + PPI & \NA & \NA & \NA & \NA & \NA & \NA \\
\bottomrule
\end{tabular}
\end{table}

Table~\ref{tab:full_grid_allocation} reports the validation-set size, the effective label budget
after removing validation labels, the LoRA-Var theoretical allocation fraction, the empirical
best allocation over the trained candidate set, the allocation actually evaluated in the performance columns, the normalized
variance ratio, allocation regret, and PPI coverage. The variance ratio is the main metric aligned
with the FT-PPI theory; values below one indicate lower estimated variance than the Sample Mean.
Regret is reported for LoRA-Var + PPI because this is the method for which the scaling-law
allocation rule is being tested. For LoRA-MSE + PPI, the row is a full-grid empirical comparator:
we report its grid-best allocation and performance there, but we do not claim a separate MSE
allocation rule.

Figure~\ref{fig:allocation_validation_b500} illustrates the full-grid validation for
$B=500$. The left panel plots the normalized estimated variance objective
$\overline V_B^{m}(\rho)/\overline V_B^{\mathrm{SM}}$ for LoRA-Var + PPI across the allocation
grid together with the exact theoretical candidate, while the right panel repeats the same display
for LoRA-MSE + PPI on the original grid. The LoRA-Var panel marks both the theoretical allocation
from \eqref{eq:target_theoretical_allocation} and the empirical best allocation over
$\mathcal{G}_B^{\mathrm{Var}}$. Analogous figures for other budgets are reported in
Appendix~\ref{app:full_grid_allocation_figures}.

\begin{figure}[H]
\centering
\begin{minipage}[t]{0.48\textwidth}
\centering
\fbox{\begin{minipage}[c][0.28\textheight][c]{0.95\linewidth}
\centering
\textbf{Panel reserved for $B=500$ LoRA-Var + PPI}\\[0.5em]
Horizontal axis: fine-tuning allocation (\%).\\
Vertical axis: normalized estimated variance $\overline V_B(\rho)/\overline V_B^{\mathrm{SM}}$.\\[0.5em]
Elements to plot: empirical variance objective, grid-best allocation,\\
and the theoretical allocation marked by a star.
\end{minipage}}

\vspace{0.4em}
(a) Variance-loss-based fine-tuning
\end{minipage}\hfill
\begin{minipage}[t]{0.48\textwidth}
\centering
\fbox{\begin{minipage}[c][0.28\textheight][c]{0.95\linewidth}
\centering
\textbf{Panel reserved for $B=500$ LoRA-MSE + PPI}\\[0.5em]
Horizontal axis: fine-tuning allocation (\%).\\
Vertical axis: normalized estimated variance $\overline V_B(\rho)/\overline V_B^{\mathrm{SM}}$.\\[0.5em]
Elements to plot: empirical variance objective and grid-best allocation.
\end{minipage}}

\vspace{0.4em}
(b) MSE-loss-based fine-tuning
\end{minipage}
\caption{Estimated FT+PPI variance objective under different sample allocations for $B=500$. The
final version will replace the two reserved panels with the corresponding full-grid objective
plots.}
\label{fig:allocation_validation_b500}
\end{figure}

\subsection{Target-Only Ramp-Up Recovery}
\label{subsec:target_rampup}

The full-grid experiment above is useful for validation, but it is not itself a deployable procedure:
it requires training at all candidate allocations before choosing one. We therefore evaluate a
target-only ramp-up procedure for LoRA-Var + PPI that estimates the allocation from the early
stages of the same target-budget split.

The ramp-up experiment uses the same label accounting as the full-grid validation. For each
budget $B$, the labeled set $\mathcal{N}_r^{(B)}$ is first split into a validation set
$\mathcal{V}_r^{(B)}$ of size $0.2B$ and an effective label set
$\mathcal{N}_{r,\mathrm{eff}}^{(B)}$ of size
\[
    n_B=B-|\mathcal{V}_r^{(B)}|=0.8B .
\]
The ramp-up training sets are nested subsets of
$\mathcal{N}_{r,\mathrm{eff}}^{(B)}$. For a candidate final allocation $s$, the fine-tuning set is
$\mathcal{N}_{r,s}^{\mathrm{FT}}$ and the PPI correction set is
$\mathcal{N}_{r,s}^{\mathrm{PPI}}
=\mathcal{N}_{r,\mathrm{eff}}^{(B)}\setminus\mathcal{N}_{r,s}^{\mathrm{FT}}$.
Thus, the validation labels are counted inside the labeled budget but are kept separate from the
correction labels.

At stage $\ell$, ramp-up has trained models only at nested sizes
$s_1<\cdots<s_\ell$ and has observed their validation residual variances on
$\mathcal{V}_r^{(B)}$. It fits
\[
    \widehat{\sigma}^2(s)=\hat a_\ell s^{-\hat\alpha_\ell}+\hat b_\ell
\]
using only these observed points, and computes the fitted minimizer of
$\widehat{\sigma}^2(s)/(B-|\mathcal{V}_r^{(B)}|-s)$. If the fitted minimizer no longer
exceeds the largest trained size, the procedure stops; otherwise it proceeds to the next
ramp-up stage.

When ramp-up stops, the final allocation is chosen from the sizes already trained by minimizing
the fitted objective over the observed grid. This best-seen rule is a conservative label-accounting
choice: labels used to train larger ramp-up models are not returned to the correction set for a
smaller final model. The validation set $\mathcal{V}_r^{(B)}$ is excluded from PPI correction to
keep allocation selection separate from inference, but it is not wasted. After the allocation is
selected, the same validation labels can continue to be used as a held-out validation set for
checkpoint selection, early stopping, or other training-control decisions for the final
fine-tuned model.

\begin{table}[H]
\centering
\caption{Target-only ramp-up recovery for LoRA-Var + PPI with validation labels counted inside
the budget}
\label{tab:rampup_recovery}
\footnotesize
\setlength{\tabcolsep}{4.5pt}
\renewcommand{\arraystretch}{1.08}
\begin{tabular}{@{}c c c c c c c c@{}}
\toprule
$B$ & $|\mathcal{V}_r^{(B)}|$ & $n_B$ & Stopping stage & Ramp-up $s$ & Oracle $s$ & Gap & Regret \\
\midrule
300 & 60 & 240 & \NA & \NA & \NA & \NA & \NA \\
500 & 100 & 400 & \NA & \NA & \NA & \NA & \NA \\
700 & 140 & 560 & \NA & \NA & \NA & \NA & \NA \\
1000 & 200 & 800 & \NA & \NA & \NA & \NA & \NA \\
\bottomrule
\end{tabular}
\end{table}

Table~\ref{tab:rampup_recovery} makes explicit that the validation set is part of the labeled
budget, not an external source of labels. The key quantity is the regret of the ramp-up-selected
allocation relative to the LoRA-Var full-grid oracle. We do not require the ramp-up estimates of
$(a,\alpha,b)$ to match the full-grid estimates exactly. What matters for inference is whether
the selected allocation achieves nearly the same objective value as the oracle allocation after
holding out the validation labels.

\subsection{End-to-End Estimator Comparison}
\label{subsec:end_to_end}

We finally evaluate the end-to-end inference performance of the resulting estimators. This
experiment follows the reporting structure of the original empirical section: we first compare
estimation accuracy, and then translate the uncertainty reduction into sample savings. The central
deployable estimator is LoRA-Var + PPI with the target-only ramp-up allocation from Section~6.4.
LoRA-MSE + PPI is included as a strong empirical comparator: for this method, we train over the
same allocation grid and report the best grid value. We do not propose an allocation rule for
LoRA-MSE + PPI, and we do not treat it as a deployable allocation method.

All methods are evaluated on the same target-population replications and budgets
\[
    B\in\{300,500,700,1000\}.
\]

\paragraph{Baselines.}
We compare against the Sample Mean, which ignores review text and uses all $B$ labeled examples
directly. We also include FT-only baselines, which average the fine-tuned model predictions over
the target unlabeled pool without PPI correction; these rows measure the bias that remains without
the correction term. The two FT+PPI rows use the same corrected estimator but differ in how the
surrogate is trained and how the fine-tuning size is selected: LoRA-MSE + PPI uses the exhaustive
grid-best value as a comparator, while LoRA-Var + PPI uses the target-only ramp-up allocation.

\paragraph{Estimation accuracy.}
For each estimator, Table~\ref{tab:estimator_accuracy} reports RMSE, bias, empirical standard
deviation, estimated standard error, and coverage. Let $\mu_{\mathrm{target}}$ denote the
finite-population mean of $\mathcal{P}_{\mathrm{target}}$ and let $\hat\mu_r$ be the estimate from
replication $r$. For example,
\[
    \mathrm{RMSE}
    =
    \sqrt{\frac{1}{R}\sum_{r=1}^R(\hat\mu_r-\mu_{\mathrm{target}})^2}.
\]

\begin{table}[H]
\centering
\caption{Estimation accuracy across estimators under different labeled budgets}
\label{tab:estimator_accuracy}
\scriptsize
\setlength{\tabcolsep}{3.5pt}
\renewcommand{\arraystretch}{1.08}
\begin{tabularx}{\textwidth}{@{}c X l c c c c c@{}}
\toprule
$B$ & Estimator & Selection & RMSE & Bias & Emp. SD & Est. SE & Coverage \\
\midrule
300 & Sample Mean & -- & \NA & \NA & \NA & \NA & \NA \\
    & LoRA-MSE FT-only & -- & \NA & \NA & \NA & \NA & \NA \\
    & LoRA-Var FT-only & -- & \NA & \NA & \NA & \NA & \NA \\
    & LoRA-MSE + PPI & grid best & \NA & \NA & \NA & \NA & \NA \\
    & LoRA-Var + PPI & ramp-up & \NA & \NA & \NA & \NA & \NA \\
\midrule
500 & Sample Mean & -- & \NA & \NA & \NA & \NA & \NA \\
    & LoRA-MSE FT-only & -- & \NA & \NA & \NA & \NA & \NA \\
    & LoRA-Var FT-only & -- & \NA & \NA & \NA & \NA & \NA \\
    & LoRA-MSE + PPI & grid best & \NA & \NA & \NA & \NA & \NA \\
    & LoRA-Var + PPI & ramp-up & \NA & \NA & \NA & \NA & \NA \\
\midrule
700 & Sample Mean & -- & \NA & \NA & \NA & \NA & \NA \\
    & LoRA-MSE FT-only & -- & \NA & \NA & \NA & \NA & \NA \\
    & LoRA-Var FT-only & -- & \NA & \NA & \NA & \NA & \NA \\
    & LoRA-MSE + PPI & grid best & \NA & \NA & \NA & \NA & \NA \\
    & LoRA-Var + PPI & ramp-up & \NA & \NA & \NA & \NA & \NA \\
\midrule
1000 & Sample Mean & -- & \NA & \NA & \NA & \NA & \NA \\
     & LoRA-MSE FT-only & -- & \NA & \NA & \NA & \NA & \NA \\
     & LoRA-Var FT-only & -- & \NA & \NA & \NA & \NA & \NA \\
     & LoRA-MSE + PPI & grid best & \NA & \NA & \NA & \NA & \NA \\
     & LoRA-Var + PPI & ramp-up & \NA & \NA & \NA & \NA & \NA \\
\bottomrule
\end{tabularx}
\vspace{0.5em}

\begin{minipage}{0.96\textwidth}
\footnotesize
Note: LoRA-MSE + PPI is reported at the best value found on the full allocation grid. It is a
comparison value, not a proposed allocation procedure. LoRA-Var + PPI uses the target-only
ramp-up allocation.
\end{minipage}
\end{table}

\paragraph{Sample savings.}
Table~\ref{tab:sample_savings} reports how much labeled data the FT+PPI estimators save relative
to the Sample Mean. We compute the savings from estimated variances: if method $m$ has estimated
variance $\widehat{\Var}(\hat\mu_m)$ and the Sample Mean under the same budget has estimated
variance $\widehat{\Var}(\hat\mu_{\mathrm{SM}})$, then
\[
    \mathrm{Savings}(m)
    =
    1-
    \frac{\widehat{\Var}(\hat\mu_m)}
         {\widehat{\Var}(\hat\mu_{\mathrm{SM}})} .
\]
Positive values indicate that the method achieves the same precision as the Sample Mean while
using fewer labeled examples. Because all methods are evaluated on matched replications, we also
report paired standard errors for the savings differences.

\begin{table}[H]
\centering
\caption{Sample savings of FT+PPI estimators relative to the Sample Mean}
\label{tab:sample_savings}
\scriptsize
\setlength{\tabcolsep}{4.5pt}
\renewcommand{\arraystretch}{1.08}
\begin{tabularx}{\textwidth}{@{}c X l c c c c@{}}
\toprule
$B$ & Estimator & Selection & CI Len. & CI ratio & Sample savings & Paired SE \\
\midrule
300 & Sample Mean & -- & \NA & 1.000 & 0.000 & \NA \\
    & LoRA-MSE + PPI & grid best & \NA & \NA & \NA & \NA \\
    & LoRA-Var + PPI & ramp-up & \NA & \NA & \NA & \NA \\
\midrule
500 & Sample Mean & -- & \NA & 1.000 & 0.000 & \NA \\
    & LoRA-MSE + PPI & grid best & \NA & \NA & \NA & \NA \\
    & LoRA-Var + PPI & ramp-up & \NA & \NA & \NA & \NA \\
\midrule
700 & Sample Mean & -- & \NA & 1.000 & 0.000 & \NA \\
    & LoRA-MSE + PPI & grid best & \NA & \NA & \NA & \NA \\
    & LoRA-Var + PPI & ramp-up & \NA & \NA & \NA & \NA \\
\midrule
1000 & Sample Mean & -- & \NA & 1.000 & 0.000 & \NA \\
     & LoRA-MSE + PPI & grid best & \NA & \NA & \NA & \NA \\
     & LoRA-Var + PPI & ramp-up & \NA & \NA & \NA & \NA \\
\bottomrule
\end{tabularx}
\end{table}

The two tables together tie the preceding experiments to the final inference task. Section~6.3
checks whether the LoRA-Var variance objective identifies a good fine-tuning fraction, and
Section~6.4 checks whether target-only ramp-up recovers a low-regret approximation to that
fraction. Tables~\ref{tab:estimator_accuracy} and \ref{tab:sample_savings} then test whether the
resulting deployable LoRA-Var + PPI estimator improves accuracy and label efficiency under the
same target-population budgets.

\subsection{Combination with Cross-PPI (To Be Determined)}
\label{subsec:crossppi_combination}

The allocation procedure above is designed for the standard FT+PPI estimator, but it is not tied
to a particular rectification implementation. A natural extension is to combine the same
fine-tuning allocation rule with Cross-PPI. This subsection is reserved for that compatibility
experiment rather than for a separate baseline comparison.

The planned experiment fixes a labeled budget, such as $B=300$, and first selects the
fine-tuning fraction using the target-only ramp-up procedure in Section~6.4. We then apply
Cross-PPI fold-wise: for each fold, the predictor used to correct that fold is trained without
using the fold's correction labels, and the corrected estimates are averaged across folds. The
large unlabeled text dataset $\mathcal{M}_r^{(B)}$ remains common across folds, while the
fold-wise role separation preserves the same no-leakage principle as the main FT+PPI estimator.

The empirical question is whether the allocation selected by our ramp-up rule remains near the
best allocation after Cross-PPI is applied. Operationally, we will plot the Cross-PPI estimator
variance over a grid of fine-tuning fractions and overlay the ramp-up-selected fraction. If, for
example, the selected fraction at $B=300$ lies near the minimum of the Cross-PPI curve, this would
show that the proposed label-allocation rule is compatible with Cross-PPI. The comparison is
therefore meant to demonstrate complementarity: Cross-PPI changes how correction labels are used,
whereas our contribution is to choose an efficient fine-tuning versus correction allocation.

\appendix
\section{Scaling-Law Validation and Diagnostic Details}
\label{app:scaling_details}

\paragraph{Training configuration.}
All models use \texttt{Qwen/Qwen2.5-1.5B-Instruct} with a scalar regression head on the last
non-padding token representation. We use LoRA adapters on all linear modules with rank 8, LoRA
scaling parameter 16, and dropout 0.05. The LoRA learning rate is $2\times10^{-4}$, the head
learning rate is $5\times10^{-4}$, weight decay is 0.01, and gradient clipping uses maximum
norm 1.0. The maximum sequence length is 192. The pilot scaling-law cells used batch size 128 on
NVIDIA L40S GPUs; no out-of-memory fallback was triggered. Training uses at most 20 epochs,
with a minimum of 4 epochs and patience 4 for early stopping.

\paragraph{Loss-function diagnostic.}
In addition to the variance-loss scaling-law run, we evaluate two MSE-based diagnostics under
the same pilot split protocol: MSE training with MSE-based checkpoint selection, and MSE
training with variance-based checkpoint selection. Table~\ref{tab:scaling_loss_fit} shows that
all three approaches improve substantially over the constant-predictor baseline, while the
variance-loss model with variance-based checkpoint selection provides the best pilot scaling fit.

\begin{table}[H]
\centering
\caption{Scaling-law fits for the pilot loss-function diagnostic}
\label{tab:scaling_loss_fit}
\small
\setlength{\tabcolsep}{7pt}
\begin{tabular}{@{}lrrrr@{}}
\toprule
Method & $\hat a$ & $\hat\alpha$ & $\hat b$ & $R^2$ \\
\midrule
MSE loss, stop by MSE & 34.69 & 0.505 & 1.993 & 0.894 \\
MSE loss, stop by Var & 19.51 & 0.491 & 2.488 & 0.989 \\
Var loss, stop by Var & 17.40 & 0.574 & 2.827 & 0.995 \\
\bottomrule
\end{tabular}
\end{table}

\paragraph{Additional diagnostics.}
\label{app:diagnostics}
The following diagnostics use the same pilot split, model family, and optimization parameters as
the scaling-law validation above. We report them together with the scaling-law details because
they check whether the fitted power law is driven by the intended residual-variance mechanism
rather than by prediction collapse, invalid rating ranges, or leakage across data roles.

First, we compare residual variance under MSE and variance fine-tuning, because residual variance
is the quantity that enters the PPI correction variance. Second, we monitor prediction variance to
detect collapse to an almost constant predictor. Third, we record the share of predictions outside
the feasible rating range $[80,100]$. Fourth, we compare every trained surrogate with the
constant predictor baseline of 9.57. Finally, we verify all role splits to rule out leakage across
pilot, validation, fine-tuning, correction, and unlabeled sets. Figure~\ref{fig:loss_diagnostic}
reports the loss-function residual-variance diagnostic, and Table~\ref{tab:diagnostics}
summarizes the remaining checks.

\begin{figure}[H]
\centering
\includegraphics[width=0.86\textwidth]{artifacts/hyak/wine_loss_controlled_b10000_r10_stopmetric_s50_to700/figures/method_comparison_validation_scale_residual_variance_raw.pdf}
\caption{Residual variance on the raw rating scale under MSE and variance fine-tuning
diagnostics. Shaded bands denote 95\% confidence intervals across 10 pilot replications.}
\label{fig:loss_diagnostic}
\end{figure}

\begin{table}[H]
\centering
\caption{Diagnostic checks for target-population experiments}
\label{tab:diagnostics}
\footnotesize
\setlength{\tabcolsep}{5pt}
\renewcommand{\arraystretch}{1.10}
\begin{tabularx}{\textwidth}{@{}p{0.20\textwidth}p{0.30\textwidth}X@{}}
\toprule
Diagnostic & Quantity & Purpose \\
\midrule
Residual variance & Var/MSE by method & Tests whether variance loss improves the PPI-relevant quantity. \\
\addlinespace[0.15em]
Prediction spread & $\widehat{\Var}\{f(X)\}$ and prediction range & Detects collapse to a nearly constant predictor. \\
\addlinespace[0.15em]
Out-of-range rate & Share of $\hat Y\notin[80,100]$ & Checks whether the regression head produces implausible ratings. \\
\addlinespace[0.15em]
Constant baseline & Residual variance relative to 9.57 & Quantifies improvement over the no-text predictor. \\
\addlinespace[0.15em]
Leakage checks & Pairwise split intersections & Verifies separation of pilot, validation, fine-tuning, correction, and unlabeled roles. \\
\bottomrule
\end{tabularx}
\end{table}

These diagnostics are necessary because valid PPI inference requires both statistical role
separation and a useful surrogate. A method that reduces prediction error without reducing
residual variance, collapses to a nearly constant predictor, or reuses validation labels during
correction would not support the intended inference claim.

\section{Additional Full-Grid Allocation Figures}
\label{app:full_grid_allocation_figures}

This appendix contains the budget-specific full-grid allocation plots not shown in the main text.
Figures~\ref{fig:allocation_validation_b300}, \ref{fig:allocation_validation_b700}, and
\ref{fig:allocation_validation_b1000} use the same visual convention as
Figure~\ref{fig:allocation_validation_b500}: the horizontal axis is the fine-tuning allocation
fraction after removing validation labels, and the vertical axis is the normalized plug-in
variance objective $\overline V_B(\rho)/\overline V_B^{\mathrm{SM}}$. The star marks the
theoretical allocation implied by the Section~6.2 scaling-law fit, and the empirical minimum over
the trained candidate set marks the full-grid benchmark.

\begin{figure}[H]
\centering
\fbox{%
\begin{minipage}[c][0.34\textheight][c]{0.82\textwidth}
\centering
\textbf{Full-grid allocation validation, $B=300$}\\[0.8em]
Plot reserved for LoRA-MSE + PPI and LoRA-Var + PPI normalized variance objectives.\\
The validation set has size 60, so the effective allocation budget is $n_B=240$.
\end{minipage}}
\caption{Planned full-grid allocation validation plot for $B=300$.}
\label{fig:allocation_validation_b300}
\end{figure}

\begin{figure}[H]
\centering
\fbox{%
\begin{minipage}[c][0.34\textheight][c]{0.82\textwidth}
\centering
\textbf{Full-grid allocation validation, $B=700$}\\[0.8em]
Plot reserved for LoRA-MSE + PPI and LoRA-Var + PPI normalized variance objectives.\\
The validation set has size 140, so the effective allocation budget is $n_B=560$.
\end{minipage}}
\caption{Planned full-grid allocation validation plot for $B=700$.}
\label{fig:allocation_validation_b700}
\end{figure}

\begin{figure}[H]
\centering
\fbox{%
\begin{minipage}[c][0.34\textheight][c]{0.82\textwidth}
\centering
\textbf{Full-grid allocation validation, $B=1000$}\\[0.8em]
Plot reserved for LoRA-MSE + PPI and LoRA-Var + PPI normalized variance objectives.\\
The validation set has size 200, so the effective allocation budget is $n_B=800$.
\end{minipage}}
\caption{Planned full-grid allocation validation plot for $B=1000$.}
\label{fig:allocation_validation_b1000}
\end{figure}

\end{document}
